\subsection{Algoritmo de Sincronización de Cristian}

\subsubsection{Lógica de Actualización de Relojes Físicos}

El algoritmo de Cristian es un método probabilístico diseñado para sincronizar relojes físicos en sistemas distribuidos. Su premisa se basa en calcular el tiempo de tránsito de un mensaje por la red, asumiendo que el viaje de ida tarda aproximadamente lo mismo que el de vuelta (simetría en la red).

Para lograr esta sincronización, el cliente registra el tiempo exacto en que envía una solicitud ($T_0$) y el tiempo en que recibe la respuesta ($T_1$). La diferencia entre ambos nos da el Tiempo de Ida y Vuelta o RTT (Round Trip Time):
\[ RTT = T_1 - T_0 \]

El servidor, al procesar la petición, inserta su propia estampa de tiempo oficial ($T_s$). Finalmente, el cliente ajusta su reloj local sumando al tiempo del servidor la mitad del RTT (que representa la latencia estimada de bajada):
\[ \text{Hora Corregida} = T_s + \frac{RTT}{2} \]

\textbf{Mecanismo de Precisión:} En el archivo \texttt{sinc\_cristian.c}, se implementó un umbral de tolerancia (\texttt{RTT\_MAX\_MS = 600}). Si el retardo supera este límite, significa que la red sufrió una congestión anómala, invalidando la asunción de simetría. En ese caso, el cliente descarta la estampa de tiempo y reintenta la operación.

\subsubsection{Arquitectura del Sistema}

La arquitectura abandona el esquema Peer-to-Peer (P2P) de los problemas anteriores para adoptar una estricta \textbf{topología en estrella} (Cliente-Servidor).
\begin{itemize}
    \item \textbf{Protocolo de Transporte:} La comunicación se realiza mediante datagramas \textbf{UDP}(\texttt{SOCK\_DGRAM}). Dado que se transmiten paquetes muy pequeños y sensibles a la latencia, UDP evita el costo temporal que implicaría establecer una conexión TCP.
    \item \textbf{Jerarquía de Procesos:} El proceso Padre asume el rol ininterrumpido de Servidor Maestro. El proceso Hijo actúa como coordinador, lanzando a los procesos nietos (quienes fungirán como clientes independientes) y, al finalizar todos, inyecta el comando \texttt{EXIT} para apagar el servidor.
\end{itemize}

\subsubsection{Estructura de Datos (\texttt{header.h})}

Para manipular el tiempo físico sin sufrir desbordamientos aritméticos, se abstrajo la hora en estructuras específicas:
\begin{itemize}
    \item \textbf{\texttt{tiempo}:} Fragmenta la hora local en horas, minutos, segundos y milisegundos, permitiendo una fácil lectura e impresión.
    \item \textbf{\texttt{Mensaje\_C}:} Es la carga útil del datagrama UDP. A diferencia de las prácticas de relojes lógicos, aquí se transportan tres estampas de tiempo vitales ($T_0$, $T_s$, $T_1$) que viajarán entre el cliente y el servidor.
\end{itemize}

\begin{lstlisting}[style=estiloC]
// ----- Estructura para almacenar el tiempo físico y visual -----
typedef struct {
    int horas;
    int minutos;
    int segundos;
    int milisegundos;
} tiempo;

// ----- Estructura del mensaje de Cristian -----------
typedef struct {
    int  id_proceso;            
    char peticion[200];         
    tiempo tiempo_envio;      // T0: Estampa de tiempo al enviar
    tiempo tiempo_servidor;   // Ts: Estampa oficial devuelta
    tiempo tiempo_respuesta;  // T1: Estampa al recibir respuesta
} Mensaje_C;
\end{lstlisting}

\subsubsection{Archivo Principal (\texttt{sinc\_cristian.c})}

El núcleo matemático del archivo consiste en la función \texttt{capturar\_tiempo}, que extrae la hora real del sistema utilizando \texttt{CLOCK\_REALTIME}, y la función \texttt{convertion\_ms}, que transforma la estructura tridimensional de tiempo a un entero de 64 bits (\texttt{long long}). Esta conversión es obligatoria para poder realizar las restas y divisiones de la fórmula de Cristian a nivel de milisegundos sin pérdida de precisión.

\begin{lstlisting}[style=estiloC]
#include <stdio.h>        // printf, perror
#include <stdlib.h>       // exit, rand, srand
#include <string.h>       // memset, strcpy, strcmp
#include <unistd.h>       // fork, sleep, usleep, getpid, close
#include <sys/wait.h>     // wait
#include <arpa/inet.h>    // inet_aton, htons
#include <netinet/in.h>   // sockaddr_in
#include <time.h>         // clock_gettime, localtime
#include "header.h"       // estructuras tiempo / Mensaje_C y constantes

// Variables globales para definir el limite en la respuesta
// Y para el retardo maximo del servidor a la hora de enviar respuestas.
#define RTT_MAX_MS  600         // 500 ms como límite razonable en loopback
#define DELAY_SERV  300000      // Retardo maximo del servidor en us
#define DELAY_CLI   250000      // Retardo maximo del cliente en us

// Esta funcion imprime una estructura tiempo con formato H:M:S:ms
void mostrar_hora(tiempo *t)
{
    printf("[%d:%d:%d:%d]\n",   // formato H:M:S:ms
           t->horas,
           t->minutos,
           t->segundos,
           t->milisegundos);
}

// Convertir estructura tiempo a milisegundos como un entero de 64 bits
long long convertion_ms(tiempo *t)
{
    // Usamos long long para numeros de 64 bits y poder tener mejor aproximacion
    // Realizamos la conversion de horas, minutos y segundos a milisegundos
    // 1 hora = 3600 segundos * 1000 ms = 3,600,000 ms
    // 1 minuto = 60 segundos * 1000 ms = 60,000 ms
    // 1 segundo = 1,000 ms
    long long t_ms = (long long)t->horas * 3600000LL +
                     (long long)t->minutos * 60000LL +
                     (long long)t->segundos * 1000LL +
                     t->milisegundos;
    return t_ms;
}

// Lee el reloj del sistema en tiempo real y llena la estructura tiempo
void capturar_tiempo(tiempo *t)
{
    struct timespec ts;                     // estructura con segundos y nanosegundos
    struct tm *tm_info;                     // estructura para desglosar hora local

    clock_gettime(CLOCK_REALTIME, &ts);     // obtiene hora del reloj del sistema
    tm_info = localtime(&ts.tv_sec);        // convierte segundos Unix a hora local

    t->milisegundos = ts.tv_nsec / 1000000; // 1 000 000 ns = 1 ms
    t->segundos     = tm_info->tm_sec;      // segundos (0-59)
    t->minutos      = tm_info->tm_min;      // minutos  (0-59)
    t->horas        = tm_info->tm_hour;     // horas    (0-23)
}

// El servidor se encarga de recibir el mensaje del cliente lo procesa y le regresa
// su respuesta con su estampa de tiempo al momento de recibir la solicitud
void ejecutar_servidor_C()
{
    int desc_socket;                 // descriptor del socket UDP del servidor
    struct sockaddr_in direct;       // dirección propia del servidor
    struct sockaddr_in cliente;      // dirección del cliente que llama
    socklen_t tam_cliente;           // tamaño de la dirección del cliente
    Mensaje_C msj;                   // buffer para el mensaje recibido/enviado

    printf("\n[SERVIDOR] Reloj Maestro iniciado en el puerto %d...\n", PUERTO);

    desc_socket = socket(AF_INET, SOCK_DGRAM, 0);   // socket UDP (sin conexión)

    memset(&direct, 0, sizeof(direct));              // limpia la estructura de dirección
    direct.sin_family      = AF_INET;                // familia IPv4
    direct.sin_port        = htons(PUERTO);          // puerto en orden de red
    direct.sin_addr.s_addr = INADDR_ANY;             // acepta peticiones de cualquier IP

    // Se asocia el socket a la direccion "escucha por el puerto"
    bind(desc_socket, (struct sockaddr *)&direct, sizeof(direct)); 

    while (1) // bucle infinito hasta recibir EXIT
    {
        tam_cliente = sizeof(cliente);  // se debe resetear el tamaño antes de cada recvfrom
        memset(&msj, 0, sizeof(Mensaje_C));  // limpia el buffer del mensaje

        // Se queda esperando hasta recibir un datagrama de algun cliente.
        recvfrom(desc_socket, &msj, sizeof(Mensaje_C), 0, (struct sockaddr *)&cliente, &tam_cliente);

        if (strcmp(msj.peticion, "EXIT") == 0) // comando de apagado remoto
        {
            printf("\n[SERVIDOR] Comando EXIT recibido. Apagando.\n");
            break; // sale del bucle y cierra el servidor
        }
        else if (strcmp(msj.peticion, "ACTUALIZAR") == 0) // petición de sincronización
        {
            // El tiempo del servidor se captura al momento de recibir la peticion
            capturar_tiempo(&msj.tiempo_servidor);   //<- Ts capturado primero

            // Simulamos latencia de procesamiento
            usleep(rand() % DELAY_SERV);
            // Mostramos sobre la terminal del servidor la hora en la que recibio la peticion con PID
            printf("\n[SERVIDOR] Petición recibida de PID %d. Hora despachada: ", msj.id_proceso);
            // Se muestra la hora [H:M:S:ms]
            mostrar_hora(&msj.tiempo_servidor);

            strcpy(msj.peticion, "OK"); // cambia la petición a confirmación

            // Devuelve el mensaje (con Ts dentro) al cliente que lo solicitó
            sendto(desc_socket, &msj, sizeof(Mensaje_C), 0, (struct sockaddr *)&cliente, tam_cliente);
        }
        // Regresamos a escuchar y limpiar la estructura.
    }
    // Si el comando EXIT se recibio, se debe cerrar el socket
    close(desc_socket); // libera el descriptor del socket
}

// El cliente se encargara de solicitar al servidor para sincronizar su reloj
void ejecutar_cliente_C(int indice)
{
    srand(time(NULL) ^ (getpid() << 16));   // semilla única por proceso
    usleep(rand() % DELAY_CLI); // retardo aleatorio para escalonar los clientes

    int desc_socket;             // descriptor del socket UDP del cliente
    struct sockaddr_in dest;     // dirección del servidor destino
    Mensaje_C msj;               // buffer del mensaje enviado/recibido
    tiempo reloj_sincronizado;   // resultado final: hora ajustada por Cristian

    desc_socket = socket(AF_INET, SOCK_DGRAM, 0);   // socket UDP

    memset(&dest, 0, sizeof(dest));         // limpia la estructura de dirección
    dest.sin_family = AF_INET;              // familia IPv4
    dest.sin_port   = htons(PUERTO);        // puerto del servidor en orden de red
    inet_aton(DIR_SERV, &dest.sin_addr);    // convierte la IP del servidor a binario

    // El algoritmo de Cristian especifica que si el RTT es demasiado grande,
    // la estimación Ts + RTT/2 tiene demasiado error y el resultado no es
    // confiable. En ese caso se debe repetir la petición hasta obtener
    // un RTT dentro del umbral aceptable (RTT_MAX_MS).
    long long rango; // RTT medido en este intento
    int intentos = 0; // contador de intentos realizados

    // Utilizamos un ciclo do-while por si la diferencia de tiempo es mayor al limite establecido
    do
    {
        intentos++; // incrementa el contador de intentos

        memset(&msj, 0, sizeof(Mensaje_C));   // limpia el mensaje antes de cada intento
        msj.id_proceso = getpid();            // identifica al proceso cliente
        strcpy(msj.peticion, "ACTUALIZAR");   // petición de sincronización al servidor

        // PASO 1: Capturar T0 y enviar la petición
        capturar_tiempo(&msj.tiempo_envio);   // T0: estampa de tiempo de salida
        sendto(desc_socket, &msj, sizeof(Mensaje_C), 0, (struct sockaddr *)&dest, sizeof(dest));

        // PASO 2: Recibir la respuesta del servidor (Ts viaja dentro de msj)
        socklen_t tam_struct = sizeof(dest);
        recvfrom(desc_socket, &msj, sizeof(Mensaje_C), 0, (struct sockaddr *)&dest, &tam_struct);

        // PASO 3: Capturamos el tiempo en que llego la respuesta
        capturar_tiempo(&msj.tiempo_respuesta);  // T1: estampa de tiempo de llegada

        // Convertimos el tiempo de envio T0 a milisegundos
        long long t0_ms = convertion_ms(&msj.tiempo_envio);

        // Convertimos el tiempo de respuesta T1 a milisegundos
        long long t1_ms = convertion_ms(&msj.tiempo_respuesta);

        rango = t1_ms - t0_ms; // RTT = T1 - T0 (diferencia de tiempo de ida y vuelta)

        if (rango > RTT_MAX_MS)  // si el rango supera el umbral muestra mensaje y repite
            printf("[CLIENTE #%02d] RTT = %lld ms demasiado alto, reintentando... (intento %d)\n", indice + 1, rango, intentos);

    } while (rango > RTT_MAX_MS); // repite hasta obtener un rango aceptable

    // FORMULA DE CRISTIAN -> si la respuesta entre el servidor y el cliente estan dentro del limite se corrige el reloj
    // El reloj del servidor mas la diferencia de tiempo ida y vuelta del mensaje entre 2 (asumiendo simetria)
    // hora_corregida = Ts + (T1 - T0) / 2
    // La mitad del RTT aproxima el tiempo que tardó el mensaje en llegar desde el servidor hasta el cliente

    // Convertimos el tiempo del servidor Ts a milisegundos
    long long ts_ms = convertion_ms(&msj.tiempo_servidor);

    long long latencia_ms      = rango / 2;            // latencia estimada = RTT / 2
    long long hora_corregida_ms = ts_ms + latencia_ms;  // hora ajustada por Cristian

    // Reconstruimos la hora corregida en la estructura tiempo (h:m:s:ms)
    reloj_sincronizado.milisegundos = hora_corregida_ms % 1000;          // residuo en ms

    long long seg_totales = hora_corregida_ms / 1000; // total en segundos
    reloj_sincronizado.segundos = seg_totales % 60;   // segundos (0-59)

    long long min_totales = seg_totales / 60;        // total en minutos
    reloj_sincronizado.minutos = min_totales % 60;   // minutos (0-59)

    long long horas_totales = min_totales / 60;    // total en horas
    reloj_sincronizado.horas = horas_totales % 24; // horas formato 24h

    // Imprimimos todos los tiempos y el resultado final
    printf("\n[CLIENTE #%02d | PID %d] ---> T0 (Salida)    : ", indice + 1, getpid());
    mostrar_hora(&msj.tiempo_envio);

    printf("[CLIENTE #%02d | PID %d] ---> T1 (Llegada)   : ", indice + 1, getpid());
    mostrar_hora(&msj.tiempo_respuesta);

    printf("[CLIENTE #%02d | PID %d] ---> Ts (Servidor)  : ", indice + 1, getpid());
    mostrar_hora(&msj.tiempo_servidor);

    printf("[CLIENTE #%02d | PID %d] ---> RTT            : %lld ms (intento(s): %d)\n", indice + 1, getpid(), rango, intentos);   // muestra el RTT final aceptado

    printf("[CLIENTE #%02d | PID %d] ---> Latencia Calc  : %lld ms\n", indice + 1, getpid(), latencia_ms);

    printf("[CLIENTE #%02d | PID %d] ---> RELOJ AJUSTADO : ", indice + 1, getpid());
    mostrar_hora(&reloj_sincronizado);

    close(desc_socket);   // libera el descriptor del socket
    usleep(1000000);      // espera 1 s para que el servidor procese antes de salir
    exit(0);              // el proceso nieto termina limpiamente
}

int main()
{
    pid_t pid = fork(); // se crea el proceso hijo (los nietos seran clientes)

    if (pid < 0) exit(1); // error al hacer fork

    // Proceso padre actuara como el servidor
    if (pid > 0)
    {
        ejecutar_servidor_C(); // Entra aquí y sale hasta recibir EXIT

        int estado;
        waitpid(pid, &estado, 0);    // espera específicamente al hijo coordinador

        printf("\n[SISTEMA] Simulación finalizada.\n");
        exit(0);
    }

    // PROCESO HIJO: coordina y lanza los clientes
    else
    {
        sleep(1);                    // espera a que el servidor esté listo

        for (int i = 0; i < N_CLIENTES; i++)       // lanza N_CLIENTES procesos nietos
        {
            if (fork() == 0)                        // cada nieto es un cliente independiente
                ejecutar_cliente_C(i);              // ejecuta el algoritmo de Cristian y hace exit(0)
        }
        // El proceso hijo espera a que todos los nietos terminen (clientes).
        for (int i = 0; i < N_CLIENTES; i++) wait(NULL);

        // Todos los clientes terminaron: ordenar el apagado del servidor
        int sock_exit = socket(AF_INET, SOCK_DGRAM, 0);  // socket temporal para enviar EXIT

        struct sockaddr_in dest_exit;
        memset(&dest_exit, 0, sizeof(dest_exit));         // limpia la dirección destino
        dest_exit.sin_family = AF_INET;                   // familia IPv4
        dest_exit.sin_port   = htons(PUERTO);             // puerto del servidor
        inet_aton(DIR_SERV, &dest_exit.sin_addr);         // IP del servidor

        Mensaje_C msj_exit;
        memset(&msj_exit, 0, sizeof(Mensaje_C));          // limpia el mensaje de salida
        strcpy(msj_exit.peticion, "EXIT");                // comando de apagado

        // Envía el EXIT al servidor para que salga de su bucle
        sendto(sock_exit, &msj_exit, sizeof(Mensaje_C), 0, (struct sockaddr *)&dest_exit, sizeof(dest_exit));

        close(sock_exit);   // libera el socket temporal
        exit(0);            // el proceso hijo coordinador termina
    }
}
\end{lstlisting}